\documentclass[10pt,a4paper]{article}
\usepackage[utf8]{inputenc}
\usepackage[T1]{fontenc}
\usepackage[left=1.8cm,right=1.8cm,top=2.06cm,bottom=1.8cm,headsep=24.3pt]{geometry}
\usepackage{xcolor,titlesec,hyperref,fancyhdr}
\definecolor{heading}{HTML}{0D4268}
\hypersetup{pdftitle={Lingua IA - HTA de correo y vocabulario v8}}
\pagestyle{fancy}\fancyhf{}
\fancyhead[L]{\small Lingua IA \textperiodcentered{} Análisis de tareas}
\fancyhead[R]{\small HTA \textperiodcentered{} Correo y vocabulario}
\fancyfoot[C]{\thepage}
\setlength{\headheight}{12pt}
\setlength{\parindent}{0pt}
\setlength{\parskip}{4pt}
\setlength{\emergencystretch}{3em}
\titleformat{\section}{\Large\bfseries\color{heading}}{}{0pt}{}
\titleformat{\subsection}{\large\bfseries\color{heading}}{}{0pt}{}
\titlespacing*{\subsection}{0pt}{10.66pt}{5.13pt}
\newcommand{\HTAline}[3]{\par\noindent\hspace*{#1cm}\begin{minipage}{\dimexpr\linewidth-#1cm\relax}#3 #2\end{minipage}\par}
\begin{document}
\hypersetup{pdfauthor={Misael Palomino Rivadeneyra, Jose Alva Cornejo}}
\textbf{Alumnos:} Misael Palomino Rivadeneyra, Jose Alva Cornejo\par
\section*{1. Análisis de potenciales grupos de usuarios (Usuarios objetivo)}
Estudiante universitario de 18 a 25 años con nivel B1-B2 de inglés.\par
\subsection*{Características generales del usuario objetivo}
a) Edad, cultura e idioma: Tiene entre 18 y 25 años y estudia en un contexto universitario hispanohablante. Su idioma principal es el español y aprende inglés como idioma adicional. Puede comprender textos generales en inglés, aunque algunas palabras, expresiones o estructuras todavía le generan dudas.\par
b) Educación: Cursa estudios superiores y tiene experiencia leyendo y escribiendo en español para fines académicos. Está familiarizado con redactar correos y textos para sus cursos. Su escritura en inglés está en desarrollo y puede variar según el tema y el tipo de texto.\par
c) Conocimientos específicos: Cuenta con conocimientos básicos e intermedios de gramática inglesa, próximos al nivel B1--B2. Puede comprender una idea general y elaborar borradores, pero necesita comprobar vocabulario contextual, combinaciones de palabras, formalidad y claridad. Debe distinguir una elección de registro de un error gramatical y de un cambio de significado.\par
d) Limitaciones físicas y mentales: La información disponible no permite afirmar que todos los estudiantes carecen de limitaciones o necesidades de accesibilidad. No se presupone que todos perciban, lean o interactúen de la misma manera.\par
e) Experiencia con la tecnología: El perfil anterior declara consultas a traductores, diccionarios y herramientas de IA para comprender vocabulario o redactar en inglés. Puede copiar textos entre aplicaciones, escribir consultas y revisar respuestas. Las rutas de esta revisión consideran esos recursos cuando están disponibles; su elección y frecuencia por caso requieren validación.\par
f) Motivación y actitud: Busca mejorar su inglés para resolver actividades académicas y desenvolverse con mayor seguridad. El alcance se centra en quien quiere comprender, comprobar y reutilizar lo obtenido. Puede trabajar con rapidez, pero recibir o entregar un texto sin revisarlo no demuestra aprendizaje. Puede frustrarse ante explicaciones ambiguas o cambios que perjudican partes que ya funcionaban.\par
g) Experiencia en el aprendizaje del inglés: Combina sus conocimientos con modelos, apuntes, fuentes de consulta y ayuda de otras personas. Puede conocer una regla sin aplicarla con seguridad. Necesita explicar sus elecciones, producir expresiones propias, recuperarlas sin el modelo y transferirlas a otra situación.\par
h) Ambiente y contexto de estudio: Puede estudiar en casa, en la universidad u otros espacios y alternar lectura, redacción y consulta de recursos. Los horarios, dispositivos y condiciones de estudio no están confirmados para todos los estudiantes. Se consideran sesiones interrumpidas y la recuperación del contexto, las notas y las dudas al retomar.\par
i) Relaciones y patrones de comunicación: Intercambia mensajes escritos con docentes, compañeros o equipos académicos. Necesita adecuar las expresiones a la confianza, la intención y los datos conocidos del destinatario. Esta caracterización toma como referencia el v7; los casos, expresiones y recorridos revisados son propuestas de análisis pendientes de contraste con usuarios.\par
\subsection*{Enfoque del análisis}
El objetivo es aprender inglés al comprender y producir correos y al estudiar vocabulario en contexto. El análisis parte de una necesidad real, un texto recibido, un borrador propio o una expresión cuyo significado o uso puede desconocerse; no presupone un correo concreto, sus palabras, sus errores ni una respuesta correcta ya identificada. Se describen acciones y decisiones reutilizables: reconocer lo que se sabe, delimitar la dificultad, recuperar o buscar información, contrastarla, intentar un uso propio y revisar lo que todavía no se comprende. Las rutas de consulta se apoyan en los recursos declarados por el perfil del v7; sus secuencias y frecuencia requieren contraste con usuarios. La explicación, recuperación y transferencia se proponen para examinar aprendizaje y no se presentan como hábitos universales observados. Los ejemplos de la interfaz ilustran estas tareas y no determinan el árbol.\par
\clearpage
\section*{Tarea 1. Aprender inglés mediante correos académicos}
Objetivo de aprendizaje: Aprender recursos lingüísticos del inglés mediante el análisis, la producción, la revisión y la reutilización de correos académicos, con atención al significado, la construcción y el registro.\par
\subsection*{Precondiciones}
{\setlength{\parskip}{2pt}\textbullet\ Una necesidad de comunicación, un correo recibido, una consigna o un borrador, sin contenido predefinido.\par}
{\setlength{\parskip}{2pt}\textbullet\ Conocimientos previos de inglés variables dentro del perfil de referencia; puede no saber formular o interpretar una expresión.\par}
{\setlength{\parskip}{2pt}\textbullet\ Posibilidad de recuperar información propia o consultar recursos disponibles. Si faltan contexto, datos o acceso a fuentes, se comprueban o registran durante la tarea.\par}
\textbf{Frecuencia:} Según necesidades académicas y sesiones de práctica; puede continuar en otra sesión.\par
\textbf{Restricciones:} El correo sitúa el aprendizaje del inglés y su envío real queda fuera del análisis. No se fijan palabras, errores, fórmulas ni destinatarios de antemano. Cada persona elige las rutas necesarias según sus conocimientos y recursos. Las secuencias son propuestas por contrastar con usuarios; completar el correo no basta como evidencia de aprendizaje.\par
\textbf{Excepciones:} Intención o destinatario inciertos; imposibilidad de iniciar una formulación; texto parcialmente comprendido; falta de referencias pertinentes; respuestas contradictorias; propuesta que altera el sentido; dificultad que continúa sin resolverse.\par
\subsection*{Análisis jerárquico de tareas}
\HTAline{0.0}{\textbf{0. Aprender inglés mediante correos académicos}}{}
\HTAline{0.0}{\textbf{1. Determinar la necesidad de aprendizaje del inglés.}}{}
\HTAline{0.4}{\textbf{1.1.} Establecer la intención comunicativa.}{}
\HTAline{0.8}{1.1.1. Definir la intención de un mensaje propio.}{}
\HTAline{0.8}{1.1.2. Interpretar la intención de un correo ajeno.}{}
\HTAline{1.2000000000000002}{1.1.2.1. Leer el correo de origen.}{}
\HTAline{1.2000000000000002}{1.1.2.2. Identificar la intención del remitente.}{}
\HTAline{1.2000000000000002}{1.1.2.3. Localizar fragmentos de comprensión incierta.}{}
\HTAline{0.8}{1.1.3. Interpretar la consigna de una actividad.}{}
\HTAline{1.2000000000000002}{1.1.3.1. Examinar la consigna de la actividad.}{}
\HTAline{1.2000000000000002}{1.1.3.2. Distinguir instrucciones explícitas y datos ausentes.}{}
\HTAline{1.2000000000000002}{1.1.3.3. Registrar ambigüedades de la consigna.}{}
\HTAline{0.4}{\textbf{1.2.} Determinar el contexto del destinatario.}{}
\HTAline{0.8}{1.2.1. Identificar la relación con el destinatario.}{}
\HTAline{0.8}{1.2.2. Distinguir datos conocidos y datos inciertos.}{}
\HTAline{0.8}{1.2.3. Comprobar los datos inciertos del contexto.}{}
\HTAline{0.4}{\textbf{1.3.} Reconocer el punto de partida lingüístico.}{}
\HTAline{0.8}{1.3.1. Recuperar una expresión de conocimientos previos.}{}
\HTAline{0.8}{1.3.2. Proponer un intento inicial en inglés.}{}
\HTAline{0.8}{1.3.3. Localizar la dificultad de significado o uso.}{}
\HTAline{0.0}{\textbf{2. Examinar referencias de uso del inglés.}}{}
\HTAline{0.4}{\textbf{2.1.} Obtener una referencia pertinente.}{}
\HTAline{0.8}{2.1.1. Localizar una referencia pertinente.}{}
\HTAline{0.8}{2.1.2. Leer la referencia de uso.}{}
\HTAline{0.8}{2.1.3. Señalar fragmentos de comprensión incierta.}{}
\HTAline{0.4}{\textbf{2.2.} Comparar formulaciones de intención equivalente.}{}
\HTAline{0.4}{\textbf{2.3.} Explicar las diferencias entre formulaciones.}{}
\HTAline{0.8}{2.3.1. Explicar diferencias de significado o registro.}{}
\HTAline{0.8}{2.3.2. Identificar límites de la explicación propia.}{}
\HTAline{0.0}{\textbf{3. Seleccionar recursos lingüísticos para el correo.}}{}
\HTAline{0.4}{\textbf{3.1.} Formular un asunto en inglés.}{}
\HTAline{0.8}{3.1.1. Precisar la intención del asunto.}{}
\HTAline{0.8}{3.1.2. Obtener recursos lingüísticos para el asunto.}{}
\HTAline{1.2000000000000002}{3.1.2.1. Recuperar una expresión de asunto conocida.}{}
\HTAline{1.2000000000000002}{3.1.2.2. Consultar alternativas para el asunto.}{}
\HTAline{1.2000000000000002}{3.1.2.3. Interpretar el significado de una alternativa.}{}
\HTAline{0.8}{3.1.3. Elaborar un asunto propio.}{}
\HTAline{1.2000000000000002}{3.1.3.1. Redactar un asunto propio.}{}
\HTAline{1.2000000000000002}{3.1.3.2. Contrastar el asunto con la intención.}{}
\HTAline{0.4}{\textbf{3.2.} Formular un saludo en inglés.}{}
\HTAline{0.8}{3.2.1. Obtener fórmulas de saludo.}{}
\HTAline{1.2000000000000002}{3.2.1.1. Recuperar una fórmula de saludo conocida.}{}
\HTAline{1.2000000000000002}{3.2.1.2. Consultar alternativas de saludo.}{}
\HTAline{1.2000000000000002}{3.2.1.3. Interpretar la función de una fórmula de saludo.}{}
\HTAline{0.8}{3.2.2. Comparar fórmulas de saludo por significado y uso.}{}
\HTAline{0.8}{3.2.3. Elaborar un saludo propio.}{}
\HTAline{1.2000000000000002}{3.2.3.1. Redactar un saludo propio.}{}
\HTAline{1.2000000000000002}{3.2.3.2. Justificar la apertura y el tratamiento.}{}
\HTAline{0.8}{3.2.4. Resolver los datos inciertos del destinatario.}{}
\HTAline{1.2000000000000002}{3.2.4.1. Comprobar los datos del destinatario.}{}
\HTAline{1.2000000000000002}{3.2.4.2. Buscar una fórmula independiente del dato incierto.}{}
\HTAline{1.2000000000000002}{3.2.4.3. Registrar la incertidumbre del destinatario.}{}
\HTAline{0.4}{\textbf{3.3.} Construir el mensaje principal en inglés.}{}
\HTAline{0.8}{3.3.1. Definir el contenido del mensaje principal.}{}
\HTAline{1.2000000000000002}{3.3.1.1. Precisar la acción o información del mensaje.}{}
\HTAline{1.2000000000000002}{3.3.1.2. Identificar los participantes del mensaje.}{}
\HTAline{1.2000000000000002}{3.3.1.3. Localizar la dificultad de formulación.}{}
\HTAline{0.8}{3.3.2. Obtener una construcción para la intención comunicativa.}{}
\HTAline{1.2000000000000002}{3.3.2.1. Recuperar una construcción conocida.}{}
\HTAline{1.2000000000000002}{3.3.2.2. Consultar construcciones de intención equivalente.}{}
\HTAline{1.2000000000000002}{3.3.2.3. Interpretar el significado y uso de una construcción.}{}
\HTAline{0.8}{3.3.3. Elaborar un intento del mensaje principal.}{}
\HTAline{1.2000000000000002}{3.3.3.1. Adaptar la construcción al contenido propio.}{}
\HTAline{1.2000000000000002}{3.3.3.2. Redactar un intento del mensaje.}{}
\HTAline{1.2000000000000002}{3.3.3.3. Anotar dudas lingüísticas del intento.}{}
\HTAline{0.8}{3.3.4. Contrastar el intento con una referencia pertinente.}{}
\HTAline{0.4}{\textbf{3.4.} Expresar relaciones de razón o finalidad.}{}
\HTAline{0.8}{3.4.1. Precisar la relación entre las ideas.}{}
\HTAline{0.8}{3.4.2. Consultar una construcción para esa relación.}{}
\HTAline{0.8}{3.4.3. Redactar una frase explicativa.}{}
\HTAline{0.8}{3.4.4. Contrastar el significado de la frase.}{}
\HTAline{0.4}{\textbf{3.5.} Formular el cierre del mensaje en inglés.}{}
\HTAline{0.8}{3.5.1. Obtener fórmulas de cierre.}{}
\HTAline{1.2000000000000002}{3.5.1.1. Identificar la función del cierre.}{}
\HTAline{1.2000000000000002}{3.5.1.2. Buscar fórmulas de cierre pertinentes.}{}
\HTAline{1.2000000000000002}{3.5.1.3. Interpretar el significado y uso de las fórmulas.}{}
\HTAline{0.8}{3.5.2. Elegir un cierre para el contexto.}{}
\HTAline{1.2000000000000002}{3.5.2.1. Comparar las alternativas de cierre.}{}
\HTAline{1.2000000000000002}{3.5.2.2. Redactar un cierre propio.}{}
\HTAline{1.2000000000000002}{3.5.2.3. Justificar la coherencia del cierre.}{}
\HTAline{0.0}{\textbf{4. Resolver una dificultad del inglés.}}{}
\HTAline{0.4}{\textbf{4.1.} Delimitar la dificultad lingüística.}{}
\HTAline{0.8}{4.1.1. Localizar el fragmento dudoso.}{}
\HTAline{0.8}{4.1.2. Conservar la oración y su contexto.}{}
\HTAline{0.8}{4.1.3. Precisar la dificultad de comprensión o producción.}{}
\HTAline{0.4}{\textbf{4.2.} Elegir un recurso de consulta.}{}
\HTAline{0.8}{4.2.1. Consultar un diccionario, apuntes o modelos.}{}
\HTAline{1.2000000000000002}{4.2.1.1. Localizar una entrada o ejemplo pertinente.}{}
\HTAline{1.2000000000000002}{4.2.1.2. Examinar el significado y la construcción.}{}
\HTAline{1.2000000000000002}{4.2.1.3. Relacionar los ejemplos con la dificultad.}{}
\HTAline{0.8}{4.2.2. Consultar un traductor.}{}
\HTAline{1.2000000000000002}{4.2.2.1. Aportar el fragmento y el contexto al traductor.}{}
\HTAline{1.2000000000000002}{4.2.2.2. Obtener una traducción del fragmento.}{}
\HTAline{1.2000000000000002}{4.2.2.3. Comparar la traducción con la intención.}{}
\HTAline{1.2000000000000002}{4.2.2.4. Señalar equivalencias inciertas.}{}
\HTAline{0.8}{4.2.3. Consultar una herramienta de IA existente.}{}
\HTAline{1.2000000000000002}{4.2.3.1. Formular una consulta con el fragmento y la dificultad.}{}
\HTAline{1.2000000000000002}{4.2.3.2. Solicitar una explicación de las alternativas.}{}
\HTAline{1.2000000000000002}{4.2.3.3. Contrastar la respuesta con el contexto o una fuente.}{}
\HTAline{0.8}{4.2.4. Consultar la ayuda de otra persona.}{}
\HTAline{1.2000000000000002}{4.2.4.1. Presentar el fragmento y la duda.}{}
\HTAline{1.2000000000000002}{4.2.4.2. Solicitar una explicación o un ejemplo.}{}
\HTAline{1.2000000000000002}{4.2.4.3. Contrastar la explicación con el contexto.}{}
\HTAline{0.4}{\textbf{4.3.} Valorar la respuesta de una consulta.}{}
\HTAline{0.8}{4.3.1. Relacionar la respuesta con el fragmento de origen.}{}
\HTAline{0.8}{4.3.2. Comprobar el respaldo de la explicación.}{}
\HTAline{0.8}{4.3.3. Explicar el hallazgo con palabras propias.}{}
\HTAline{0.8}{4.3.4. Identificar dudas persistentes.}{}
\HTAline{0.4}{\textbf{4.4.} Reformular una consulta insuficiente.}{}
\HTAline{0.8}{4.4.1. Precisar el contexto faltante.}{}
\HTAline{0.8}{4.4.2. Solicitar una explicación alternativa.}{}
\HTAline{0.8}{4.4.3. Contrastar otra fuente.}{}
\HTAline{0.8}{4.4.4. Registrar la duda persistente.}{}
\HTAline{0.0}{\textbf{5. Revisar una producción propia en inglés.}}{}
\HTAline{0.4}{\textbf{5.1.} Preparar el borrador para la revisión.}{}
\HTAline{0.8}{5.1.1. Redactar un borrador propio.}{}
\HTAline{0.8}{5.1.2. Conservar la versión original del borrador.}{}
\HTAline{0.8}{5.1.3. Identificar fragmentos de revisión.}{}
\HTAline{0.4}{\textbf{5.2.} Examinar las diferencias lingüísticas.}{}
\HTAline{0.8}{5.2.1. Contrastar una construcción gramatical.}{}
\HTAline{1.2000000000000002}{5.2.1.1. Recuperar una regla gramatical conocida.}{}
\HTAline{1.2000000000000002}{5.2.1.2. Consultar la regla gramatical pertinente.}{}
\HTAline{1.2000000000000002}{5.2.1.3. Comparar la regla con el fragmento.}{}
\HTAline{0.8}{5.2.2. Evaluar la pertinencia del registro.}{}
\HTAline{1.2000000000000002}{5.2.2.1. Comparar el efecto de las alternativas.}{}
\HTAline{1.2000000000000002}{5.2.2.2. Distinguir una elección de uso de un error gramatical.}{}
\HTAline{0.8}{5.2.3. Comprobar la fidelidad del significado.}{}
\HTAline{1.2000000000000002}{5.2.3.1. Comparar participantes, acciones, condiciones y datos.}{}
\HTAline{1.2000000000000002}{5.2.3.2. Señalar alteraciones de la intención.}{}
\HTAline{0.4}{\textbf{5.3.} Decidir sobre el fragmento de revisión.}{}
\HTAline{0.8}{5.3.1. Incorporar una alternativa pertinente.}{}
\HTAline{0.8}{5.3.2. Conservar la formulación propia.}{}
\HTAline{0.8}{5.3.3. Solicitar una aclaración de la propuesta.}{}
\HTAline{0.4}{\textbf{5.4.} Documentar una decisión lingüística.}{}
\HTAline{0.8}{5.4.1. Explicar la razón de la decisión lingüística.}{}
\HTAline{0.8}{5.4.2. Conservar la comparación de versiones.}{}
\HTAline{0.0}{\textbf{6. Reutilizar recursos del inglés en otra situación.}}{}
\HTAline{0.4}{\textbf{6.1.} Recuperar una expresión sin modelo.}{}
\HTAline{0.8}{6.1.1. Apartar las referencias de apoyo.}{}
\HTAline{0.8}{6.1.2. Formular un intento de memoria.}{}
\HTAline{0.8}{6.1.3. Identificar lagunas de recuperación.}{}
\HTAline{0.4}{\textbf{6.2.} Explicar el uso de la expresión del intento.}{}
\HTAline{0.8}{6.2.1. Explicar el significado y la adecuación del intento.}{}
\HTAline{0.8}{6.2.2. Registrar las ayudas del intento.}{}
\HTAline{0.4}{\textbf{6.3.} Transferir un recurso a otra situación.}{}
\HTAline{0.8}{6.3.1. Elegir una situación de transferencia.}{}
\HTAline{0.8}{6.3.2. Adaptar la expresión al nuevo contexto.}{}
\HTAline{0.8}{6.3.3. Producir un mensaje propio.}{}
\HTAline{0.4}{\textbf{6.4.} Contrastar los usos del recurso.}{}
\HTAline{0.8}{6.4.1. Comparar el significado entre situaciones.}{}
\HTAline{0.8}{6.4.2. Explicar la adaptación del recurso.}{}
\HTAline{0.8}{6.4.3. Identificar dudas de transferencia.}{}
\HTAline{0.0}{\textbf{7. Retomar el aprendizaje del inglés.}}{}
\HTAline{0.4}{\textbf{7.1.} Conservar la intención, relación, versión, explicación y dudas de la práctica.}{}
\HTAline{0.4}{\textbf{7.2.} Priorizar una duda de la sesión anterior.}{}
\HTAline{0.8}{7.2.1. Identificar aprendizajes previos.}{}
\HTAline{0.8}{7.2.2. Seleccionar una dificultad de continuidad.}{}
\HTAline{0.4}{\textbf{7.3.} Recuperar una expresión en otra sesión.}{}
\HTAline{0.8}{7.3.1. Formular una expresión de memoria.}{}
\HTAline{0.8}{7.3.2. Registrar la fecha y las ayudas del intento.}{}
\HTAline{0.4}{\textbf{7.4.} Valorar la continuidad del aprendizaje.}{}
\HTAline{0.8}{7.4.1. Contrastar el intento con referencias pertinentes.}{}
\HTAline{0.8}{7.4.2. Elegir la siguiente práctica.}{}
\subsection*{Planes de aprendizaje}
\textbf{Plan 0:} Realizar 1 para establecer el punto de partida. Si necesita comprender o encontrar expresiones, realizar 2 y las partes pertinentes de 3; si ya dispone de un borrador, comenzar su contraste en 5 y volver a 2 o 3 donde falte comprensión. Intercalar 4 cuando desconozca una expresión, una construcción o un criterio de elección, y volver al punto de origen. Intentar 6 para recuperar y reutilizar lo comprendido; conservar y retomar lo pendiente mediante 7. No se presupone un contenido, una corrección ni una secuencia idénticos para todos los casos.\par
\textbf{Plan 1:} Realizar 1.1 y 1.2; en 1.3 reconocer el punto de partida lingüístico. Si no puede formular nada en inglés, expresar la intención con sus recursos habituales y acudir a 2 o 4; no exigir un intento inventado. Si faltan datos de la situación, buscarlos o mantenerlos pendientes. Elegir la siguiente tarea según necesite comprender un texto, encontrar una expresión o revisar un borrador.\par
\textbf{Plan 1.1:} Realizar 1.1.1, 1.1.2 o 1.1.3 según el origen de la necesidad; combinar las rutas cuando el caso lo requiera. Si no puede reconocer la intención del texto, delimitar en 4 qué impide comprenderla y volver a esta tarea. Si la consigna sigue siendo ambigua, comprobarla con la fuente o conservar la duda, sin completar por suposición lo que se espera comunicar.\par
\textbf{Plan 1.1.2:} Realizar 1.1.2.1, 1.1.2.2 y 1.1.2.3. Si la intención no se reconoce, acudir a 4 con el fragmento incierto y volver a 1.1.2.2; no completar el mensaje por suposición.\par
\textbf{Plan 1.1.3:} Realizar 1.1.3.1, 1.1.3.2 y 1.1.3.3. Comprobar las ambigüedades con la fuente cuando sea posible; conservar como pendiente lo que no tenga confirmación.\par
\textbf{Plan 1.2:} Realizar 1.2.1 y 1.2.2. Realizar 1.2.3 solo ante información incierta; si no se obtiene confirmación, conservarla como pendiente en 7.1. No suponer nombres, títulos o grados de confianza.\par
\textbf{Plan 1.3:} Realizar 1.3.1 si dispone de una expresión conocida o 1.3.2 si puede elaborar un intento. Continuar con 1.3.3. Si no puede comenzar en inglés, expresar la intención con sus recursos habituales y acudir a 2 o 4; no exigir una producción inventada.\par
\textbf{Plan 2:} Realizar 2.1. Si necesita elegir entre formulaciones, buscar una alternativa pertinente y realizar 2.2; comprobar mediante 2.3 qué puede explicar de la comparación. Si solo encuentra una referencia, estudiar su significado y uso sin inventar otra. Acudir a 4 ante información no comprendida o fuentes contradictorias. Si el recurso ya se conoce y resulta pertinente, continuar en 3 o 5 sin comparar por obligación.\par
\textbf{Plan 2.1:} Realizar 2.1.1, 2.1.2 y 2.1.3. Elegir una referencia de intención y relación comparables. Si no hay ninguna disponible o no se comprende, acudir a 4; no asumir que existe un modelo preparado para el caso.\par
\textbf{Plan 2.3:} Realizar 2.3.1 y 2.3.2 sobre las alternativas de 2.2. Si la explicación carece de respaldo o solo existe una referencia, consultar 4 sin inventar otra formulación.\par
\textbf{Plan 3:} Elegir las partes necesarias: 3.1 si no sabe expresar el asunto; 3.2 si necesita un saludo; 3.3 para el mensaje principal; 3.4 si requiere una razón o finalidad; 3.5 si necesita un cierre. Se pueden abordar en el orden en que aparezcan las dificultades y revisar unas después de otras. Intercalar 4 si desconoce el significado o uso de una opción. Reunir lo comprendido en 5; no se exige estudiar todas las partes ni conocer de antemano las fórmulas que aparecerán.\par
\textbf{Plan 3.1:} Realizar 3.1.1. En 3.1.2 recuperar una expresión que conoce o consultar si no dispone de ella; no suponer que el vocabulario ya está identificado en inglés. Continuar con 3.1.3 cuando pueda explicar lo que propone. Si no entiende la referencia o la adaptación cambia el sentido, volver a 4 o 3.1.1. Si el asunto existente ya expresa la intención y no plantea una dificultad, conservarlo.\par
\textbf{Plan 3.1.2:} Elegir 3.1.2.1 si dispone de una expresión pertinente; de lo contrario realizar 3.1.2.2. Realizar 3.1.2.3 sobre la opción disponible, con ayuda de 4 si fuera necesaria. No asumir palabras de antemano.\par
\textbf{Plan 3.1.3:} Realizar 3.1.3.1 y 3.1.3.2. Si la formulación cambia la intención o incorpora información sin respaldo, volver a 3.1.1 o 3.1.2.\par
\textbf{Plan 3.2:} Realizar 3.2.1 y 3.2.2 sobre las alternativas disponibles. Si faltan datos del destinatario o de la relación, realizar 3.2.4 antes de decidir. Continuar con 3.2.3 cuando pueda justificar una opción. Si solo conoce una fórmula y no sabe si encaja, consultar 4 en lugar de atribuirle un registro. Si el dato o el uso siguen sin confirmarse, mantener la elección pendiente.\par
\textbf{Plan 3.2.1:} Elegir 3.2.1.1 si dispone de una fórmula pertinente o 3.2.1.2 si necesita una referencia. Realizar 3.2.1.3 antes de la comparación; consultar 4 si no comprende el uso.\par
\textbf{Plan 3.2.3:} Realizar 3.2.3.1 y 3.2.3.2. Si la justificación depende de datos desconocidos, volver a 3.2.4; si depende de un significado incierto, acudir a 4.\par
\textbf{Plan 3.2.4:} Realizar 3.2.4.1. Si no hay confirmación, intentar 3.2.4.2 sin inventar información. Realizar 3.2.4.3 si el dato continúa incierto; suspender la elección de tratamiento si depende de ese dato.\par
\textbf{Plan 3.3:} Realizar 3.3.1 y elegir en 3.3.2 entre recuperar una construcción conocida o consultar cómo expresar la intención. Si no comprende la estructura, realizar 4 antes de adaptarla. Elaborar 3.3.3 y contrastar mediante 3.3.4; si el significado o la construcción no se pueden justificar, volver a 3.3.1 o 4. Repetir solo sobre la dificultad pendiente; no sustituir por defecto todo el mensaje.\par
\textbf{Plan 3.3.1:} Realizar 3.3.1.1, 3.3.1.2 y 3.3.1.3 en ese orden, sobre el mismo fragmento o contexto.\par
\textbf{Plan 3.3.2:} Elegir 3.3.2.1 si dispone de una construcción pertinente o 3.3.2.2 si la desconoce. Realizar 3.3.2.3 sobre la opción disponible; si la explicación no se comprende, acudir a 4 antes de la adaptación.\par
\textbf{Plan 3.3.3:} Realizar 3.3.3.1, 3.3.3.2 y 3.3.3.3 en ese orden, sobre el mismo fragmento o contexto.\par
\textbf{Plan 3.4:} Realizar 3.4.1 cuando la intención requiera una razón o finalidad. Realizar 3.4.2 solo si desconoce una construcción pertinente; acudir a 4 si no comprende la referencia. Continuar con 3.4.3 y 3.4.4. Si la relación cambia, volver a 3.4.1. Omitir esta rama cuando no sea necesaria.\par
\textbf{Plan 3.5:} Realizar 3.5.1 y después 3.5.2. Si desconoce la función o el uso de una fórmula, consultar 4 antes de elegir. Si las referencias no permiten distinguir opciones o contradicen la relación, ampliar el contexto en 1.2 o conservar la duda. No añadir una despedida diferente solo porque existe otra alternativa.\par
\textbf{Plan 3.5.1:} Realizar 3.5.1.1. Si no dispone de una fórmula pertinente, realizar 3.5.1.2; continuar con 3.5.1.3 sobre las fórmulas disponibles. Consultar 4 ante un significado o uso desconocidos.\par
\textbf{Plan 3.5.2:} Realizar 3.5.2.1 cuando existan alternativas. Continuar con 3.5.2.2 y 3.5.2.3. Si solo existe una referencia, justificar su pertinencia sin inventar otra. Volver a 4 si el uso sigue incierto.\par
\textbf{Plan 4:} Realizar 4.1. Si la dificultad exige determinar el significado o uso de vocabulario, acudir a raíz 0 de la tarea 2 con la oración y el punto de origen; al obtener una interpretación o una duda explícita, continuar en 4.3. Para otras dificultades elegir una vía de 4.2 y contrastar en 4.3. Si no se comprende la respuesta, realizar 4.4 y repetir la consulta o cambiar de recurso. Volver a la tarea de origen solo con lo que pueda justificar; conservar en 7 lo que no se resuelva.\par
\textbf{Plan 4.1:} Realizar 4.1.1, 4.1.2 y 4.1.3. Si no puede clasificar la dificultad como vocabulario, gramática o registro, conservar su descripción abierta antes de 4.2; no imponer una categoría.\par
\textbf{Plan 4.2:} Elegir una vía disponible de 4.2.1--4.2.4 según la dificultad y los recursos habituales; no son pasos acumulativos obligatorios. Si desconoce qué buscar, aportar la oración y describir qué no comprende. Si la respuesta no permite comprobar el uso, consultar otra vía o precisar la pregunta. Si no hay acceso a una fuente suficiente, registrar el pendiente en 7. La disponibilidad y preferencia por estas vías deben contrastarse con usuarios.\par
\textbf{Plan 4.2.1:} Realizar 4.2.1.1, 4.2.1.2 y 4.2.1.3 en la fuente disponible. Si la dificultad exige comparar acepciones, acudir a raíz 0 de la tarea 2 con la oración de origen. Si no existe una entrada pertinente, elegir otra ruta de 4.2.\par
\textbf{Plan 4.2.2:} Realizar 4.2.2.1, 4.2.2.2, 4.2.2.3 y 4.2.2.4. Si el nivel de inglés no permite comprobar una equivalencia, conservar la duda y contrastar otra vía de 4.2; no considerar una traducción aislada como explicación suficiente.\par
\textbf{Plan 4.2.3:} Realizar 4.2.3.1 y 4.2.3.2 cuando esta herramienta forme parte de los recursos disponibles. Realizar 4.2.3.3 antes de la incorporación de una respuesta. Si la explicación es insuficiente o contradictoria, acudir a 4.4.\par
\textbf{Plan 4.2.4:} Realizar 4.2.4.1, 4.2.4.2 y 4.2.4.3 si esa ayuda está disponible. Si no se comprende la explicación o persiste una duda, acudir a 4.4 o conservar el pendiente en 7.1.\par
\textbf{Plan 4.3:} Realizar 4.3.1, 4.3.2, 4.3.3 y 4.3.4. Si el respaldo falta o no se comprende, continuar en 4.4. Si hay una explicación suficiente, volver a la tarea de origen con el hallazgo y sus límites.\par
\textbf{Plan 4.4:} Realizar 4.4.1 cuando falte contexto. Elegir 4.4.2 si la explicación es insuficiente o 4.4.3 si necesita otro respaldo; se pueden combinar ambas rutas. Repetir 4.2 y 4.3 cuando haya nueva información. Realizar 4.4.4 si no hay una resolución suficiente, con cierre parcial en 7.1.\par
\textbf{Plan 5:} Si ya existe un borrador, conservarlo en 5.1; si no existe, producirlo con lo comprendido en 1--3 y la ayuda necesaria de 4. Revisar los fragmentos dudosos mediante 5.2, decidir en 5.3 y explicar en 5.4. Si no reconoce una dificultad, contrastar con una referencia o solicitar revisión sin asumir un error concreto. Volver a 2, 3 o 4 si no comprende una propuesta. Pasar a 6 con el recurso que pueda explicar; terminar el correo no resuelve por sí solo el aprendizaje pendiente.\par
\textbf{Plan 5.1:} Realizar 5.1.1 si no existe un borrador; si ya existe, comenzar en 5.1.2. Continuar con 5.1.3 según las dudas del estudiante o una revisión disponible, sin asumir un error específico. Acudir a 2, 3 o 4 ante dificultades de redacción.\par
\textbf{Plan 5.2:} Comparar el fragmento y la alternativa. Aplicar 5.2.1 para una posible diferencia gramatical, 5.2.2 para una elección de uso y 5.2.3 para comprobar que se conserva el significado; pueden concurrir varias. Si no puede determinar qué cambió, consultar 4 sin clasificarlo a la fuerza. No presuponer qué errores contiene el borrador ni dar por correcta toda propuesta recibida.\par
\textbf{Plan 5.2.1:} Elegir 5.2.1.1 cuando conozca la regla o 5.2.1.2 cuando necesite una explicación. Continuar con 5.2.1.3. Si la regla o su aplicación no se comprenden, acudir a 4 y conservar la clasificación como incierta.\par
\textbf{Plan 5.2.2:} Realizar 5.2.2.1 y 5.2.2.2 para una misma relación e intención. Si no puede justificar la diferencia, consultar 4; no considerar toda alternativa más cortés como una corrección obligatoria.\par
\textbf{Plan 5.2.3:} Realizar 5.2.3.1 y 5.2.3.2. Si una propuesta altera el sentido o incorpora datos sin respaldo, conservar la formulación anterior o solicitar aclaración mediante 5.3.\par
\textbf{Plan 5.3:} Elegir 5.3.1 o 5.3.2 después de 5.2 según la pertinencia de la propuesta. Si falta comprensión, realizar 5.3.3 y volver a 4 o 5.2. Comparar con la versión previa; si la alternativa perjudica el sentido, recuperar la formulación anterior. La decisión debe poder explicarse en 5.4.\par
\textbf{Plan 5.4:} Realizar 5.4.1 y 5.4.2 en ese orden, sobre el mismo fragmento o contexto.\par
\textbf{Plan 6:} Intentar 6.1 y explicar el resultado en 6.2, identificando cualquier consulta. Realizar 6.3 cuando haya un recurso comprendido que pueda utilizar en otra situación; comprobar el nuevo uso mediante 6.4. Si no recuerda o no puede adaptarlo, volver a 2, 3 o 4 y conservar la dificultad. La recuperación y transferencia se proponen como actividades de aprendizaje, no como hábitos ya observados en todos los usuarios.\par
\textbf{Plan 6.1:} Realizar 6.1.1, 6.1.2 y 6.1.3. Si necesita ayuda, consultar la referencia pertinente y señalarla en 6.2.2. Un intento con ayuda conserva ese carácter; no se supone independiente por el cierre de esta tarea.\par
\textbf{Plan 6.2:} Realizar 6.2.1 y 6.2.2 en ese orden, sobre el mismo fragmento o contexto.\par
\textbf{Plan 6.3:} Realizar 6.3.1, 6.3.2 y 6.3.3 cuando haya un recurso comprendido y pertinente. Si el nuevo contexto exige otra construcción o acepción, volver a 2, 3 o 4; no forzar el recurso.\par
\textbf{Plan 6.4:} Realizar 6.4.1, 6.4.2 y 6.4.3 en ese orden, sobre el mismo fragmento o contexto.\par
\textbf{Plan 7:} Al terminar realizar 7.1. En otro momento realizar 7.2 y 7.3; elegir 4 si persiste una duda, 6 si se necesita otra recuperación o 1 si cambia la situación. Realizar 7.4 después de contrastar el intento. Una actividad interrumpida conserva su punto de retorno; no se marca dominada por cerrar la sesión.\par
\textbf{Plan 7.2:} Realizar 7.2.1 y 7.2.2 a partir de la información de 7.1. Priorizar una duda de significado, construcción o registro; si cambia la necesidad comunicativa, volver a 1.\par
\textbf{Plan 7.3:} Realizar 7.3.1 antes de la relectura del modelo; continuar con 7.3.2. Se permite ayuda y debe quedar identificada. Si la recuperación no resulta suficiente, volver a 4 o 6 sin reiniciar todo el análisis.\par
\textbf{Plan 7.4:} Realizar 7.4.1 y 7.4.2 según significado, adecuación y necesidad de ayuda. El cierre de una sesión o un intento aislado no certifican dominio; conservar cualquier pendiente mediante 7.1.\par
\subsection*{Evidencias de aprendizaje}
Necesidad y punto de partida identificados; dificultad expresada; interpretación o intento propio; recurso de consulta utilizado; comparación justificada de significado, construcción o registro; decisión de mantener o reformular un fragmento; duda restante y, cuando se realice, intento de recuperación o reutilización con la ayuda identificada.\par
Intención y dificultad identificadas; fragmento propio o interpretación con una explicación respaldada; dudas y ayuda utilizadas registradas. Los intentos de recuperación o reutilización permiten comprobar qué sigue necesitando consulta. Un correo completado, una traducción recibida o una propuesta aceptada no bastan para afirmar aprendizaje.\par
\clearpage
\section*{Tarea 2. Aprender vocabulario inglés en contexto académico}
Objetivo de aprendizaje: Aprender vocabulario inglés en contexto académico mediante inferencia, contraste, producción propia y recuperación, con atención al significado, al patrón de uso y a la ayuda necesaria.\par
\subsection*{Precondiciones}
{\setlength{\parskip}{2pt}\textbullet\ Una expresión desconocida o de uso incierto encontrada en un texto o necesaria para una producción propia.\par}
{\setlength{\parskip}{2pt}\textbullet\ Un fragmento inicial y una necesidad de comprensión o uso; no se exige conocer ya el tema, la función gramatical o la acepción.\par}
{\setlength{\parskip}{2pt}\textbullet\ Posibilidad de ampliar contexto, consultar un recurso o mantener la duda pendiente si no hay información suficiente.\par}
\textbf{Frecuencia:} Durante lecturas o redacción, con intentos posteriores de recuperación.\par
\textbf{Restricciones:} No se fija una palabra ni su significado antes de analizar el contexto. Una traducción aislada o una primera acepción no confirman el uso. Se puede declarar que no se pudo inferir o resolver la duda. Los patrones se contrastan con fuentes y la reutilización con otro contexto.\par
\textbf{Excepciones:} Contexto insuficiente; varias acepciones compatibles; ausencia de etiqueta de registro; traducción aislada; patrón no respaldado; palabra forzada en otro contexto; uso dependiente del modelo.\par
\subsection*{Análisis jerárquico de tareas}
\HTAline{0.0}{\textbf{0. Aprender vocabulario inglés en contexto académico}}{}
\HTAline{0.0}{\textbf{1. Contextualizar una dificultad de vocabulario.}}{}
\HTAline{0.4}{\textbf{1.1.} Interpretar el fragmento de lectura.}{}
\HTAline{0.8}{1.1.1. Leer el fragmento disponible.}{}
\HTAline{0.8}{1.1.2. Identificar el tema y la situación.}{}
\HTAline{0.8}{1.1.3. Señalar información contextual faltante.}{}
\HTAline{0.4}{\textbf{1.2.} Priorizar una expresión relevante.}{}
\HTAline{0.4}{\textbf{1.3.} Delimitar la duda de vocabulario.}{}
\HTAline{0.8}{1.3.1. Conservar la oración de origen.}{}
\HTAline{0.8}{1.3.2. Describir la dificultad de significado o uso.}{}
\HTAline{0.0}{\textbf{2. Inferir un significado provisional.}}{}
\HTAline{0.4}{\textbf{2.1.} Examinar la función contextual de la expresión.}{}
\HTAline{0.8}{2.1.1. Observar la posición de la expresión.}{}
\HTAline{0.8}{2.1.2. Relacionar la expresión con palabras próximas.}{}
\HTAline{0.8}{2.1.3. Proponer una función gramatical.}{}
\HTAline{0.4}{\textbf{2.2.} Localizar pistas de significado.}{}
\HTAline{0.8}{2.2.1. Relacionar la expresión con acciones y referentes de la oración.}{}
\HTAline{0.8}{2.2.2. Reconocer reformulaciones o contrastes del fragmento.}{}
\HTAline{0.8}{2.2.3. Relacionar componentes o combinaciones con conocimientos previos.}{}
\HTAline{0.4}{\textbf{2.3.} Formular una hipótesis de significado.}{}
\HTAline{0.8}{2.3.1. Proponer un significado provisional.}{}
\HTAline{0.8}{2.3.2. Justificar la hipótesis con una pista.}{}
\HTAline{0.8}{2.3.3. Registrar límites de la hipótesis.}{}
\HTAline{0.0}{\textbf{3. Contrastar la interpretación del vocabulario.}}{}
\HTAline{0.4}{\textbf{3.1.} Elegir un recurso de contraste.}{}
\HTAline{0.8}{3.1.1. Consultar las acepciones y ejemplos del diccionario.}{}
\HTAline{0.8}{3.1.2. Comparar la traducción con la oración de origen.}{}
\HTAline{0.8}{3.1.3. Consultar acepciones mediante una IA existente.}{}
\HTAline{1.2000000000000002}{3.1.3.1. Aportar la oración y el contexto.}{}
\HTAline{1.2000000000000002}{3.1.3.2. Solicitar diferencias entre acepciones.}{}
\HTAline{0.8}{3.1.4. Solicitar una explicación léxica de otra persona.}{}
\HTAline{0.4}{\textbf{3.2.} Relacionar las acepciones con el contexto.}{}
\HTAline{0.8}{3.2.1. Comparar las acepciones con la oración y los ejemplos.}{}
\HTAline{0.8}{3.2.2. Conservar la referencia de consulta.}{}
\HTAline{0.4}{\textbf{3.3.} Determinar la interpretación contextual.}{}
\HTAline{0.8}{3.3.1. Justificar una interpretación con respaldo contextual.}{}
\HTAline{0.8}{3.3.2. Diferenciar interpretaciones compatibles.}{}
\HTAline{1.2000000000000002}{3.3.2.1. Comparar diferencias entre acepciones.}{}
\HTAline{1.2000000000000002}{3.3.2.2. Localizar pistas discriminantes.}{}
\HTAline{0.8}{3.3.3. Investigar la falta de respaldo contextual.}{}
\HTAline{1.2000000000000002}{3.3.3.1. Ampliar el fragmento de origen.}{}
\HTAline{1.2000000000000002}{3.3.3.2. Consultar una explicación alternativa.}{}
\HTAline{1.2000000000000002}{3.3.3.3. Registrar el significado pendiente.}{}
\HTAline{0.4}{\textbf{3.4.} Reformular la interpretación del vocabulario.}{}
\HTAline{0.8}{3.4.1. Comparar la hipótesis con la interpretación contextual.}{}
\HTAline{0.8}{3.4.2. Explicar el significado con palabras propias.}{}
\HTAline{0.8}{3.4.3. Conservar incertidumbres restantes.}{}
\HTAline{0.0}{\textbf{4. Caracterizar el uso del vocabulario.}}{}
\HTAline{0.4}{\textbf{4.1.} Comparar usos de una misma acepción.}{}
\HTAline{0.4}{\textbf{4.2.} Determinar el registro con respaldo de la fuente.}{}
\HTAline{0.8}{4.2.1. Relacionar la etiqueta de uso con la acepción y el ejemplo.}{}
\HTAline{0.8}{4.2.2. Describir la situación de uso sin etiqueta de registro.}{}
\HTAline{0.4}{\textbf{4.3.} Deducir un patrón de uso contextual.}{}
\HTAline{0.8}{4.3.1. Localizar combinaciones de la expresión.}{}
\HTAline{0.8}{4.3.2. Examinar la construcción de los ejemplos.}{}
\HTAline{0.8}{4.3.3. Contrastar coincidencias entre los usos.}{}
\HTAline{0.8}{4.3.4. Formular un patrón provisional.}{}
\HTAline{0.4}{\textbf{4.4.} Documentar el aprendizaje léxico.}{}
\HTAline{0.8}{4.4.1. Redactar una nota de significado y uso.}{}
\HTAline{0.8}{4.4.2. Conservar el contexto y la fuente.}{}
\HTAline{0.8}{4.4.3. Registrar dudas de la expresión.}{}
\HTAline{0.0}{\textbf{5. Practicar vocabulario en una producción propia.}}{}
\HTAline{0.4}{\textbf{5.1.} Elaborar una oración propia.}{}
\HTAline{0.8}{5.1.1. Elegir una idea para la expresión.}{}
\HTAline{0.8}{5.1.2. Formular una oración de significado pertinente.}{}
\HTAline{0.8}{5.1.3. Señalar necesidades de apoyo.}{}
\HTAline{0.4}{\textbf{5.2.} Contrastar la oración con ejemplos de la misma acepción.}{}
\HTAline{0.4}{\textbf{5.3.} Decidir la formulación de la oración.}{}
\HTAline{0.8}{5.3.1. Mantener una formulación pertinente.}{}
\HTAline{1.2000000000000002}{5.3.1.1. Conservar la oración propia.}{}
\HTAline{1.2000000000000002}{5.3.1.2. Justificar una coincidencia con el ejemplo.}{}
\HTAline{0.8}{5.3.2. Reformular la construcción de la oración.}{}
\HTAline{0.8}{5.3.3. Sustituir una expresión inadecuada.}{}
\HTAline{0.0}{\textbf{6. Transferir vocabulario a otro contexto.}}{}
\HTAline{0.4}{\textbf{6.1.} Recuperar una expresión sin definición.}{}
\HTAline{0.8}{6.1.1. Apartar la definición y los apuntes.}{}
\HTAline{0.8}{6.1.2. Explicar el significado de memoria.}{}
\HTAline{0.8}{6.1.3. Producir una frase de memoria.}{}
\HTAline{0.4}{\textbf{6.2.} Definir el contexto de transferencia.}{}
\HTAline{0.8}{6.2.1. Elegir una situación diferente.}{}
\HTAline{0.8}{6.2.2. Precisar el significado de la nueva producción.}{}
\HTAline{0.4}{\textbf{6.3.} Aplicar la expresión en un nuevo contexto.}{}
\HTAline{0.8}{6.3.1. Producir un texto con la expresión.}{}
\HTAline{0.8}{6.3.2. Explicar el significado contextual.}{}
\HTAline{0.8}{6.3.3. Sustituir una expresión improcedente.}{}
\HTAline{0.8}{6.3.4. Consultar otra acepción.}{}
\HTAline{0.4}{\textbf{6.4.} Contrastar el resultado de transferencia.}{}
\HTAline{0.8}{6.4.1. Comparar el nuevo uso con el contexto original.}{}
\HTAline{0.8}{6.4.2. Registrar el contenido de memoria.}{}
\HTAline{0.8}{6.4.3. Registrar las ayudas del intento.}{}
\HTAline{0.0}{\textbf{7. Retomar el aprendizaje del vocabulario.}}{}
\HTAline{0.4}{\textbf{7.1.} Conservar el contexto del aprendizaje léxico.}{}
\HTAline{0.8}{7.1.1. Guardar la palabra, la oración, el significado y la producción.}{}
\HTAline{0.8}{7.1.2. Registrar dudas y estado de la práctica.}{}
\HTAline{0.4}{\textbf{7.2.} Recuperar el contexto de la sesión anterior.}{}
\HTAline{0.8}{7.2.1. Recordar el propósito de la práctica anterior.}{}
\HTAline{0.8}{7.2.2. Evocar la expresión desde una pista.}{}
\HTAline{0.4}{\textbf{7.3.} Practicar el recuerdo del vocabulario.}{}
\HTAline{0.8}{7.3.1. Explicar el significado de memoria.}{}
\HTAline{0.8}{7.3.2. Producir una oración de memoria.}{}
\HTAline{0.8}{7.3.3. Anotar la fecha y las ayudas del intento.}{}
\HTAline{0.4}{\textbf{7.4.} Decidir la siguiente práctica de vocabulario.}{}
\subsection*{Planes de aprendizaje}
\textbf{Plan 0:} Realizar 1 y 2 para situar la oración y reconocer el punto de partida; continuar con 3 y 4 para contrastar significado y uso. Si la consulta directa es necesaria, pasar de 1 a 3 sin inventar una hipótesis. Continuar con 5 y 6 para producción, recuperación y transferencia; cerrar o retomar mediante 7. Si se inició desde 4 de la tarea 1 y hay que volver al correo después de 4, devolver una explicación o una duda explícita y dejar el resto pendiente. No declarar terminada la raíz por haber consultado una definición.\par
\textbf{Plan 1:} Realizar 1.1 para reconocer lo que comprende y 1.2 para priorizar una dificultad; conservar el fragmento y la duda en 1.3. Si no reconoce siquiera el tema, ampliar el contexto, buscar una referencia comprensible o pedir ayuda antes de atribuir un significado. Si no hay una palabra relevante para la necesidad actual, dejarla anotada y continuar. Si procede de 4.1 de la tarea 1, conservar el punto de retorno al correo.\par
\textbf{Plan 1.1:} Realizar 1.1.1, 1.1.2 y 1.1.3. Si no se reconoce siquiera el tema, ampliar el contexto o solicitar ayuda mediante 3.1 antes de atribuir un significado.\par
\textbf{Plan 1.3:} Realizar 1.3.1 y 1.3.2. Si la consulta procede de 4.1 de la tarea 1, conservar también ese punto de retorno. Si falta contexto, ampliar el fragmento en 1.1 antes de la elección de una acepción.\par
\textbf{Plan 2:} Realizar 2.1 y 2.2 antes de 2.3. Mantener separadas hipótesis y explicación consultada. Si la persona ya conoce una acepción o decide consultar directamente, registrar ese punto de partida y pasar a 3; no presentar una hipótesis reconstruida después como si fuera inicial.\par
\textbf{Plan 2.1:} Realizar 2.1.1 y 2.1.2; intentar 2.1.3 con esas pistas. Si no puede reconocer la función, mantenerla como incierta para 3 en lugar de imponer una categoría.\par
\textbf{Plan 2.2:} Aplicar 2.2.1, 2.2.2 o 2.2.3 según las pistas presentes; se pueden combinar. La pista debe proceder de la oración, el contexto o un conocimiento previo pertinente, sin asumirla como prueba definitiva. Justificar su relación con la hipótesis en 2.3. Si no hay respaldo suficiente, ampliar 1.3 o continuar en 3 con la incertidumbre.\par
\textbf{Plan 2.3:} Intentar 2.3.1 y 2.3.2 cuando existan pistas. Realizar 2.3.3 si la evidencia es limitada; si no hay una inferencia posible, continuar en 3 con esa incertidumbre, sin inventar una hipótesis.\par
\textbf{Plan 3:} Realizar 3.1, 3.2 y 3.3 antes de 3.4. Si varias interpretaciones siguen siendo compatibles, obtener más contexto y repetir el contraste. Si no hay evidencia suficiente, registrar la duda en 7 en vez de afirmar una acepción. Pasar a 4 cuando se pueda explicar la relación entre el significado elegido y su oración.\par
\textbf{Plan 3.1:} Elegir una ruta disponible de 3.1.1--3.1.4. Registrar de dónde procede la explicación; no llamar diccionario a una respuesta generada. Si solo se obtiene una traducción aislada o una afirmación sin apoyo contextual, pedir un ejemplo o contrastar con otra vía antes de 3.3.\par
\textbf{Plan 3.1.3:} Realizar 3.1.3.1 y 3.1.3.2 si la herramienta está disponible. Contrastar en 3.2 y 3.3 la respuesta obtenida; no denominar diccionario a una respuesta generada ni asumir su exactitud.\par
\textbf{Plan 3.2:} Realizar 3.2.1 y 3.2.2 en ese orden, sobre el mismo fragmento o contexto.\par
\textbf{Plan 3.3:} Elegir 3.3.1 cuando una interpretación tenga apoyo suficiente; realizar 3.3.2 si varias siguen siendo compatibles y 3.3.3 si ninguna se puede justificar. En los dos últimos casos volver a 1.3 o 3.1 con la información que falta y repetir el contraste. Si la incertidumbre continúa, registrarla en 7 sin afirmar una acepción. No se presupone la palabra, el conjunto de acepciones ni la respuesta correcta.\par
\textbf{Plan 3.3.2:} Realizar 3.3.2.1 y 3.3.2.2 cuando varias acepciones sigan siendo posibles. Si no hay una pista suficiente, ampliar contexto en 1.3 o consultar 3.1 antes de una nueva decisión.\par
\textbf{Plan 3.3.3:} Realizar 3.3.3.1 si falta contexto o 3.3.3.2 si necesita otra explicación; se pueden combinar ambas. Volver a 3.3 con la información nueva. Realizar 3.3.3.3 si la incertidumbre persiste, con continuidad en 7.\par
\textbf{Plan 3.4:} Realizar 3.4.1 si existe una hipótesis previa; no reconstruirla después de la consulta. Continuar con 3.4.2 y 3.4.3. Volver a 3.1 o 3.3 si no hay respaldo suficiente.\par
\textbf{Plan 4:} Realizar 4.1 y 4.2 sobre la acepción elegida; continuar con 4.3 y 4.4. Si los ejemplos no sostienen un patrón, buscar otro en 3 antes de generalizar. Si se necesita volver al correo de inmediato, devolver la explicación de 4.4 a 4.3 de la tarea 1 y dejar 5--7 pendientes de práctica; no dar por completado 0. Comparar ejemplos de la misma acepción; si la diferencia exige otro significado, volver a 3. Un cambio de acepción no demuestra por sí solo una diferencia de registro.\par
\textbf{Plan 4.2:} Elegir 4.2.1 únicamente cuando exista esa información en la fuente; en caso contrario realizar 4.2.2. Si dos fuentes discrepan, conservar sus referencias y la duda. No deducir formalidad del tamaño de la palabra ni confundir la cortesía de toda la oración con una etiqueta de la palabra aislada.\par
\textbf{Plan 4.3:} Realizar 4.3.1, 4.3.2 y 4.3.3 sobre ejemplos de la misma acepción. Intentar 4.3.4 solo con respaldo suficiente; consultar 3 ante discrepancias o conservar la generalización como pendiente.\par
\textbf{Plan 4.4:} Realizar 4.4.1, 4.4.2 y 4.4.3 en ese orden, sobre el mismo fragmento o contexto.\par
\textbf{Plan 5:} Realizar 5.1, 5.2 y 5.3. Si no se puede producir sin copiar, usar un ejemplo como apoyo y declararlo; la recuperación sin definición se intentará después en 6. Pasar a 6 cuando haya una producción propia y una explicación de al menos una decisión de uso.\par
\textbf{Plan 5.1:} Realizar 5.1.1, 5.1.2 y 5.1.3. Si la expresión no corresponde a la idea, elegir otra en 5.3.3. Si necesita un modelo, identificar ese apoyo en lugar de presentar una copia como recuperación independiente.\par
\textbf{Plan 5.3:} Elegir 5.3.1, 5.3.2 o 5.3.3 después de 5.2. Si no se comprende el motivo de la revisión, volver a 3 o 4. Tras reformular, repetir 5.2 sobre la nueva frase. Conservar la producción previa para explicar qué cambió.\par
\textbf{Plan 5.3.1:} Realizar 5.3.1.1 y 5.3.1.2 en ese orden, sobre el mismo fragmento o contexto.\par
\textbf{Plan 6:} Intentar 6.1 apartando los apoyos; si necesita consultarlos, identificarlo. Precisar otra situación mediante 6.2, producir 6.3 y contrastar el resultado en 6.4. Si la nueva situación exige otra acepción o patrón, volver a 3--5; no forzar la expresión. Si el texto propio es un correo que requiere revisar la relación o el registro, acudir a 1.2 de la tarea 1 y 5.2 de la tarea 1. La práctica posterior se propone para comprobar reutilización, no se presume realizada por todos los usuarios.\par
\textbf{Plan 6.1:} Intentar 6.1.1, 6.1.2 y 6.1.3. Si necesita ayuda, consultar la referencia pertinente y registrar esa ayuda en 6.4.3. La falta de recuerdo orienta una nueva práctica y no se sustituye por una respuesta inventada.\par
\textbf{Plan 6.2:} Realizar 6.2.1 y 6.2.2 en ese orden, sobre el mismo fragmento o contexto.\par
\textbf{Plan 6.3:} Realizar 6.3.1 y 6.3.2. Si la expresión no comunica la intención, realizar 6.3.3 o 6.3.4 según se necesite otra palabra o una acepción distinta; en este último caso volver a 3. No forzar el uso de la palabra de la práctica anterior.\par
\textbf{Plan 6.4:} Realizar 6.4.1, 6.4.2 y 6.4.3 en ese orden, sobre el mismo fragmento o contexto.\par
\textbf{Plan 7:} Al terminar o interrumpir realizar 7.1. En otro momento realizar 7.2 y 7.3, después 7.4. Si la actividad quedó tras 4, retomar 5 antes de valorar la recuperación. Volver a 3 ante dudas de acepción, a 4 ante dudas de patrón o a 6 si se necesita otro intento. Guardar una palabra o abrir una nota no demuestra aprendizaje; el intervalo de repaso no se considera eficaz sin evaluar los intentos.\par
\textbf{Plan 7.1:} Realizar 7.1.1 y 7.1.2 en ese orden, sobre el mismo fragmento o contexto.\par
\textbf{Plan 7.2:} Realizar 7.2.1 y 7.2.2 sin mostrar primero la definición. Si la actividad quedó parcial, retomar el punto pendiente de 3, 4 o 5 antes de una comprobación de recuerdo.\par
\textbf{Plan 7.3:} Intentar 7.3.1 y 7.3.2 antes del contraste con los apuntes; continuar con 7.3.3. Se permite apoyo y debe quedar identificado. Pasar a 7.4 para la decisión de continuidad.\par
\subsection*{Evidencias de aprendizaje}
Hipótesis justificada o incertidumbre inicial explícita; acepción contrastada; diferencia entre significado y registro; patrón con ejemplos; anotación propia; oración producida y revisada; transferencia a otro contexto y resultado de un repaso posterior con la ayuda identificada.\par
Interpretación explicada con una pista, contraste con fuente o ejemplo, producción propia y estado de ayuda o pendiente. Guardar la palabra no demuestra que ya pueda recuperarse.\par
\end{document}